\section{Abstract Model Classes, Characterization, and Definability}
\label{sec:p03-lindstrom-continuation}

Isomorphism invariance, definability in a logic, and invariance under changes
of the ambient set universe are different requirements. A definition which
works on models of one fixed cardinality need not give a single definition
at all cardinalities. The model-class and cardinal-dependent constructions
below use the presentation of
\cite[\S2, Theorems~2 and~7; \S5, Theorem~18]{KennedyVaananen2021LogicalityV1}.
Their \(\Sigma\)- and \(\Delta\)-operations are semantic projections,
not first-order existential quantifiers over relations.

\subsection{Logic hypotheses and comparison of presentations}

\begin{hbconvention}{def:p03-lindstrom-standing}{Structures, signatures, and the ambient metatheory}
Work in \(\mathsf{ZF}+\mathsf{AC}\). Structures have nonempty domains
and set-sized finitary relational vocabularies. The explicit invariant-class
coding first uses the source's binary vocabulary \(\{R\}\); the
finite-relational generalization is stated separately. A model class \(K\)
is a class of structures in one vocabulary closed under isomorphism.
Class predicates may have fixed set parameters; restricting to a set of
structure codes uses Separation. No disjunction indexed by a proper class
is admitted as a formula.
In the native displays, \(M\) names a structure, while the same letter
in \(M^k\), \(|M|\), and the carrier component of a binary-relation
pair denotes its underlying domain. Those carrier occurrences do not
discard or replace the structure's relation interpretations.
\end{hbconvention}

\begin{hbdefinition}{def:p03-lindstrom-source-logic}{The source pair and the regular relational package}
The source writes a logic as \(L^*=(\Sigma,T)\): \(\Sigma\) is a set
or class of sentences and \(T\) relates a sentence to a structure.
Here \(T\) is a \emph{satisfaction relation}, not a theory of axioms.
Its model classes are \(\operatorname{Mod}(\phi)=\{M:T(\phi,M)\}\).
The source's five class-level conditions are isomorphism closure,
conjunction, negation, vocabulary permutations and free expansion.

For the compactness/LS characterization use the stronger regular relational
package of \cref{thm:p03-lindstrom}: sentences have finite occurrence
vocabularies; truth depends on that reduct; arity-preserving renaming
commutes with truth; relational first-order logic is included; negation and
conjunction are available; fresh constants can be abstracted with the usual
quantifier semantics; and unary relativization evaluates a sentence on the
nonempty induced domain of its guard. Countable compactness means
simultaneous satisfiability of every countable finitely satisfiable set of
sentences. Sentence-level countable LS means that every satisfiable sentence
has a countable model. Free expansion alone is not unary relativization.
These conditions are hypotheses on syntax and satisfaction, not a definition
which already assumes first-order expressiveness.
\end{hbdefinition}

\begin{hblemma}{thm:p03-lindstrom-class-bridge}{proved=annex:p03-lindstrom-class-bridge}{The class-level comparison and its extra hypotheses}
For a fixed vocabulary \(\tau\), set
\(\Sigma=\mathcal L^*[\tau]\) and
\(T(\phi,M)\) iff \(M\models_{\mathcal L^*}\phi\).
The regular package supplies the five source class-level operations with
their exact model-class meanings. This does not make the five conditions
alone imply finite occurrence, relativization, countable compactness or LS.
\end{hblemma}

\begin{hbtheorem}{thm:p03-lindstrom-generalized-class}{proved=annex:p03-lindstrom-generalized-class}{A class quantifier is not a compactness construction}
For binary vocabulary, an isomorphism-closed class \(K\) is definable
in an extension of first-order logic by the generalized quantifier
\(Q_K\), with the source semantics
\[
 M,s\models Q_Kxy\,\phi(x,y,\bar a)
 \quad\Longleftrightarrow\quad
 (M,\{(b,c)\in M^2:M,s[x\mapsto b,y\mapsto c]\models\phi\})\in K.
\]
In particular \(Q_Kxy\,R(x,y)\) defines \(K\).
Conversely every class defined in an isomorphism-invariant logic is closed
under isomorphism. This establishes source Theorem~2's equivalence of
invariance and definability in \emph{some} logic; it supplies neither
compactness nor an LS theorem for that extension.
\end{hbtheorem}

\subsection{Actual ordinal and relation codes}

\begin{hbdefinition}{def:p03-lindstrom-infinitary}{Set-sized infinitary formulas}
In \(L_{\kappa\lambda}\), conjunctions and disjunctions have fewer than
\(\kappa\) entries and quantifier blocks have fewer than \(\lambda\)
variables. \(L_{\infty\omega}\) allows arbitrary set-sized conjunctions
but finite quantifier blocks, with finitely many free variables per formula.
\(L_{\infty\infty}\) also allows set-sized variable blocks. Formula
trees are well founded; bound variables are renamed before substitution.
We use empty disjunction as falsity and empty conjunction as truth.
These syntax classes need not have finite vocabulary support.
\end{hbdefinition}

\begin{hbdefinition}{def:p03-lindstrom-ordinal}{The source ordinal formulas with an explicit order guard}
Let \(\operatorname{LO}(<)\) be the first-order axioms for a strict linear
order on the whole domain. The source's recursive formulas (2) and (4) are
\[
\begin{split}
 \eta_\alpha(x)&\leftrightarrow
 \forall y\left(y<x\to\bigvee_{\beta<\alpha}\eta_\beta(y)\right)
       \land\bigwedge_{\beta<\alpha}\exists y(y<x\land\eta_\beta(y)),\\
 \eta^0_\alpha&\leftrightarrow
 \forall y\bigvee_{\beta<\alpha}\eta_\beta(y)
       \land\bigwedge_{\beta<\alpha}\exists y\eta_\beta(y).
\end{split}
\]
The displays specify formulas by well-founded transfinite recursion; their
\(\leftrightarrow\) is not an extra logical primitive. When used on an
arbitrary expansion the order guard is included explicitly.
\end{hbdefinition}

\begin{hbtheorem}{thm:p03-lindstrom-ordinal-code}{proved=annex:p03-lindstrom-ordinal-code}{The ordinal code is an actual order-type test}
In a strict linear order, \(\eta_\alpha(a)\) holds exactly when the
predecessor segment of \(a\) has order type \(\alpha\).
The sentence \(\operatorname{LO}(<)\land\eta^0_\alpha\) holds exactly
on orders isomorphic to \(\alpha\). There is at most one point of any
specified predecessor order type. The order hypothesis cannot be dropped:
a two-point antichain satisfies \(\eta^0_1\) but is not the ordinal one.
\end{hbtheorem}

\begin{hbdefinition}{def:p03-lindstrom-cardinal-dependent}{Cardinal dependence and semantic projections}
For a cardinal \(\lambda\), the source writes
\(K_\lambda=\{M\in K:|M|=\lambda\}\).
CD-definability permits a different defining sentence for each cardinal,
not a proper-class disjunction of those sentences. A class is
\(\Sigma(L^*)\)-definable if it is the projection, by forgetting added
symbols, of an \(L^*\)-definable class of expansions on the same domain.
It is \(\Delta(L^*)\)-definable if both it and its complement are
\(\Sigma(L^*)\)-definable. Complements below are taken in all nonempty
structures of the base vocabulary. A cardinal-relative pair which only
partitions models of size \(\lambda\) is distinguished explicitly from
such a global complement pair.
\end{hbdefinition}

\begin{hbtheorem}{thm:p03-lindstrom-mcgee-code}{proved=annex:p03-lindstrom-mcgee-code}{A complete fixed-cardinality isomorphism sentence}
For an infinite cardinal \(\lambda\), every binary-relation structure
\(A\) of size \(\lambda\) has a sentence
\(\theta_A\in L_{\lambda^+\lambda^+}\) characterizing its isomorphism
type. For any isomorphism-closed \(K\), a disjunction of these sentences
over the set of its codes on domain \(\lambda\) defines \(K_\lambda\)
in \(L_{(2^\lambda)^+\lambda^+}\).
The sentence explicitly supplies distinct enumerating variables, a domain
cover clause, and the entire positive/negative relation diagram.
\end{hbtheorem}

\begin{hbtheorem}{thm:p03-lindstrom-delta-code}{proved=annex:p03-lindstrom-delta-code}{Ordinal expansions give both projection witnesses}
For each infinite cardinal \(\lambda\), the class \(K_\lambda\)
of an isomorphism-closed binary model class is
\(\Delta(L_{\infty\omega})\)-definable. On the size-\(\lambda\)
models the source's \(\Phi_A,\Theta_K\) and ordinal order code give
the positive and negative projection witnesses. The annex also supplies
the additional smaller/larger-cardinality branches needed to make the
negative projection the complement in \emph{all} base structures.
The finite-relational version follows by the explicitly proved arity bridge.
\end{hbtheorem}

\subsection{Set-theoretic definition versus sentence admission}

\begin{hbtheorem}{thm:p03-lindstrom-absolute-truth}{proved=annex:p03-lindstrom-absolute-truth}{Finite assignments give absolute infinitary truth}
For the stated relational \(L_{\infty\omega}\) syntax, well-formed
formula admission has a \(\Sigma_1\) definition in set theory and
satisfaction has a \(\Delta_1\) definition over \(\mathsf{ZF}\),
with the formula, structure and finite assignment as parameters.
Thus the source's Definition~17 absoluteness applies at this declared
set-theoretic regime. If two transitive ZF universes contain the same
admitted formula, structure and assignment, they agree on its truth.
The proof uses finite assignments at every node, not a domain of all
external infinite sequences.
\end{hbtheorem}

\begin{hbtheorem}{thm:p03-lindstrom-delta-absoluteness}{proved=annex:p03-lindstrom-delta-absoluteness}{A fixed complement pair and its admission distinction}
Fix two \(L_{\infty\omega}\) expansion sentences whose projections
are disjoint and exhaustive on the base model class in question.
Membership and nonmembership have \(\Sigma_1\) expansion-witness
definitions. For transitive universes in which the fixed pair is admitted
as a complementary pair, membership is absolute. This semantic assertion
does not make the predicate admitting \emph{all} such pairs
\(\Sigma_1\): exhaustivity ranges over all base structures, not just
over the syntax tree. In particular cardinal-relative witnesses cannot
be used to assert unconditional all-cardinality or all-universe admission
\cite[Definition~17 and the discussion preceding Theorem~18, pp.~26--27]
{KennedyVaananen2021LogicalityV1}.
\end{hbtheorem}

\subsection{Compactness and actual elementary amalgamation}

\begin{hbtheorem}{thm:p03-lindstrom-amalgamation}{proved=annex:p03-lindstrom-amalgamation}{What the regular characterization transfers to arbitrary theories}
Assume the full regular relational package, countable compactness and
sentence-level countable LS stated above. Then arbitrary set-sized
\(\mathcal L^*\)-theories satisfy compactness. A map is elementary
for \(\mathcal L^*\) exactly when it is first-order elementary, testing
sentences with each finite named parameter tuple. If
\(C\preccurlyeq_{\mathcal L^*}A,B\) in the same vocabulary, there is
a structure \(D\) and \(\mathcal L^*\)-elementary embeddings of
\(A,B\) into \(D\) agreeing on \(C\). For countable vocabulary and
countable input models, \(D\) may be countable.
The supplied diagram construction is binary elementary amalgamation;
it does not assert the entire theory of compact logics from those words alone.
\end{hbtheorem}

\subsection{Worked codes and assumption failures}

\begin{hbtheorem}{thm:p03-lindstrom-assumptions}{proved=annex:p03-lindstrom-assumptions}{The infinity quantifier retains LS but destroys compactness}
Let \(Q_0x\phi\) mean that at least \(\aleph_0\) members satisfy
\(\phi\), as in the source's \(Q_\alpha\) convention. The logic
\(L(Q_0)\) has the regular relational operations and the countable
sentence-level LS property in countable vocabulary, but not countable
compactness. The witness-closure construction of the countable model
is supplied in the annex. It strictly extends first-order expressiveness.
The sentence \(Q_1x(x=x)\), with the existing uncountability semantics,
independently demonstrates failure of countable LS; no full compactness
proof for \(L(Q_1)\) is asserted here.
\end{hbtheorem}

\begin{hbtheorem}{thm:p03-lindstrom-two-point-code}{proved=annex:p03-lindstrom-two-point-code}{A complete binary relation code and a failed elementary amalgam premise}
For the two-point structure with relation
\(R=\{(0,0),(0,1)\}\), the ordinal expansion uses the four signed
atoms \(R(x,y),R(x,y),\neg R(x,y),\neg R(x,y)\), respectively
at position pairs \((0,0),(0,1),(1,0),(1,1)\). It codes exactly that
isomorphism type, and the positive/negative projection pair is explicit.
For amalgamation, mere common substructure is insufficient: a one-point
empty unary-predicate structure is a substructure both of itself and of
a two-point structure with a new predicate member, but the two elementary
diagrams over that common substructure are incompatible.
\end{hbtheorem}

\begin{proofboundary}{def:p03-lindstrom-regimes}{Uniformity and signature boundaries}
These cardinal-dependent formulas are not a uniform all-cardinality
definition. A generalized class quantifier is not an existence proof for
compactness or completeness. The ordinal expansion has a real linear-order
guard; finite syntax support is a different condition from finite quantifier
blocks. The regular compactness/LS package is stronger than the five
class-level operations. These distinctions prevent an invariant class,
or a named semantic projection, from standing in for a characterization
construction.
\end{proofboundary}
